\documentclass[11pt,a4paper]{article}
\usepackage[T1]{fontenc}
\usepackage{microtype}
\usepackage{isabelle,isabellesym}
\usepackage{amsmath}
\usepackage{amssymb}
\usepackage{url}
\usepackage{pdfsetup}

\urlstyle{rm}
\isabellestyle{it}
% Keep long Isabelle identifiers readable in the checked source listing.
\def\isacharunderscore{\discretionary{-}{}{-}}
\def\isacharunderscorekeyword{\discretionary{-}{}{-}}
\emergencystretch=3em

\begin{document}

\title{Verified Cached Inference for Decoder Transformers}
\author{Arthur Freitas Ramos \and David Barros Hulak \and
  Ruy J. G. B. de Queiroz}
\maketitle

\begin{abstract}
This entry formalizes exact-real inference semantics for decoder-only
transformers and proves that incremental key--value caching refines
full-prefix evaluation.  The refinement is lifted from causal attention
through dimension-explicit multi-head attention, decoder blocks, and finite
stacks.  It covers rotary grouped-query attention and GPT-Neo with global or
sliding-window context, including projected per-head cache storage and finite
greedy generation.  The modern path includes an explicit arbitrary-N-step
theorem equating full-prefix first-argmax generation with the cached run under
its generation-state invariant.  Separate numerical theorems bound selected
dyadic and IEEE projection traces.
\end{abstract}

\tableofcontents

\section{Overview}

Decoder-only transformers generate tokens from left to right
\cite{VaswaniEtAl2017}.  At a new
position, causal attention needs the current query and the keys and values
already produced by the prefix.  Caching those projections is an optimization
only when it preserves the full-prefix result.  The principal theorem here
states that preservation under an explicit cache invariant and also proves the
corresponding cache-update invariant.

\section{Formalization}

The generic cache lemmas are instantiated with finite-list tensors, exact-real
scaled dot-product softmax, normalization, feed-forward layers, residual
blocks, and finite stacks.  The modern path models grouped-query attention
and rotary position embeddings~\cite{AinslieEtAl2023,SuEtAl2024}, RMS
normalization, and SwiGLU~\cite{Shazeer2020}.

The GPT-Neo path models learned token and position embeddings, affine
LayerNorm, GELU-New, biased ordinary multi-head attention, and global or
sliding-window context.  Its implementation-level cache is a list of
per-head projected key--value histories; an earlier normalized-input cache is
retained as a semantic bridge for the refinement proof.  For a positive
window, tail trimming is append-stable and idempotent.  The resulting
theorems cover attention, blocks, stacks, prompt evaluation, logits, and
finite greedy generation.  Both the modern and GPT-Neo paths expose N-step
full-versus-cached first-argmax refinements.  Synthetic checkpoints exercise
the integration.

\section{Related work}

Brucker and Stell formalize feed-forward neural networks in Isabelle/HOL,
including equivalent graph- and layer-based models and TensorFlow.js model
import~\cite{BruckerStell2023,NeuralNetworksAFP2025}.  That development
provides general neural-network representations and verification properties;
the present entry concentrates on autoregressive sequence semantics and the
refinement of a concrete inference optimization.

TorchLean gives Lean~4 a broader shared vocabulary for typed tensors, graph
intermediate representations, execution, finite-precision semantics, and
certificate checking~\cite{GeorgeEtAl2026TorchLean}.  The two efforts are
complementary: this Isabelle development isolates the list-level theorem that
projected per-head KV storage and bounded tail retention preserve decoder
semantics, while its numerical results deliberately expose their finite
precision boundary.

\section{Numerical scope}

The exact semantic theorems are complemented by numerical refinement results.
Column L1 norms give a Lipschitz constant for the vocabulary projection.  A
fixed-fraction dyadic fused-multiply-add kernel contributes at most
$d/2^{p+1}$ per logit, while the AFP IEEE model gives a $d\epsilon$ bound when
finite-value, range, and nearest-value witness conditions are supplied
\cite{YuIEEE2013}.  Replayable certificates discharge those local obligations
for selected projection traces.

The public TinyStories-1M artifact~\cite{EldanLi2023TinyStoriesCheckpoint} is
deliberately scoped: it certifies four selected layer-zero Q/K/V/MLP binary32
rows with pinned provenance and replayable FMA witnesses.  It is not a
complete attention, GELU, softmax, or all-layer floating-point replay.  The
synthetic fixtures provide the complete semantic cache and generation
integration independently of that limitation.

\section{Trusted boundary}

The formal claims have three deliberately separated trust boundaries.  First,
the exact-real cache and generation theorems are Isabelle/HOL definitions and
kernel-checked proofs over mathematical real vectors and matrices.  In
particular, they prove semantic equality between full-prefix and cached
execution; they do not assert that a particular CUDA, CPU, or PyTorch kernel
has the same floating-point evaluation order.

Second, checkpoint extraction is an external, reproducible boundary.
\texttt{tools/import\_frozen\_checkpoint.py} validates the pinned configuration,
checkpoint SHA-256 digest, tensor shapes, binary32 dtype, finite-value
preconditions, and emits exact sign/exponent/fraction triples.  Isabelle then
checks the generated data-bearing theory and its FMA certificate.  The
importer and the hash comparison are not Isabelle theorems: they establish
which bytes were extracted, while the kernel checks the mathematical
consequences of the emitted fields.

Third, the IEEE results use the AFP floating-point model's bounded-format
\texttt{fmul\_add RNE} together with explicit finite, range, and nearest-value
witnesses.  They are conditional certificate theorems, not a proof of
hardware conformance or a repaired tie-breaking implementation.  The
TinyStories result is therefore labelled a replayable local FMA certificate
demonstration.  Its bound $64$ is the conservative sum of one unit budget for
each of 64 certified steps, not an observed numerical-accuracy estimate or an
end-to-end GPT-Neo error bound.

\section{Checked Isabelle source}

The remainder of this document is the Isabelle-generated source of the checked
session.  It contains the definitions, lemmas, and proofs whose kernel-checked
build establishes the results summarized above.

\sloppy
\input{session}
\fussy

\nocite{Shazeer2019}
\bibliographystyle{abbrv}
\bibliography{root}

\end{document}
